% Part: proof-theory
% Chapter: sequent-calculus
% Section: introduction

\documentclass[../../../include/open-logic-section]{subfiles}

\begin{document}

\olfileid{pt}{seq}{int}

\olsection{جان پچھان}

سیکوئنٹ حساب !!{proof} نظام نیں جیہڑے گیرہارٹ گینٹسن نے 1930 دیاں دہائیاں وچ متعارف کرائے۔ ایہناں دا مقصد !!{proving} لئی اک عملی طریقہ دینا گھٹ سی۔ گینٹسن نے ایہ ایس لئی ورتے کہ !!{proof}s دی ممکن بناوٹ تے اوہناں دی خاصیت بارے کجھ نتیجے ثابت کر سکے۔ عین ایہو جیہے سوال ثبوت دے نظریے دا موضوع نیں۔

جیویں ناں توں پتا لگدا اے، سیکوئنٹ حساب \emph{سیکوئنٹاں} اُتے کم کردے نیں، !!{formula}s اُتے نہیں۔ سیکوئنٹ !!{formula}s دیاں لڑیاں یا کثیر سیٹاں دا جوڑا اے۔ گینٹسن دے اصل نظام (\Log{LK} تے \Log{LJ}) لڑیاں ورتدے سن۔ ہُن ثبوت دے نظریے دے ماہر کثیر سیٹ ورتنا پسند کردے نیں۔ کثیر سیٹ، سیٹ ورگا اے، فرق ایہ اے کہ کثیر سیٹ کسے !!{element} نوں اک توں ودھ واری رکھ سکدا اے۔ پر سیٹاں وانگ !!{element}s دی ترتیب نال فرق نہیں پیندا۔ سو عبارتیں $\{!A, !A, !B\}$ تے $\{!A, !B, !A\}$ اوہی کثیر سیٹ دَسّدیاں نیں، پر اوہ کثیر سیٹ $\{!A,
!B\}$ تے $\{!A, !A, !B, !B\}$ توں وکھرا اے، کیونکہ پچھلے دوہاں وچ $!A$ تے $!B$ دیاں گنتیاں پہلے نالوں وکھریاں نیں۔

رواج مطابق ثبوت دے نظریے وچ !!{formula}s دے کثیر سیٹ وڈّے یونانی حروف نال لکھے جاندے نیں۔ سو جدوں اسیں $\Gamma$، $\Delta$، $\Pi$، $\Lambda$ یا $\Theta$ لکھدے آں، تاں اوہناں توں ساڈا مطلب اوس منطق دے !!{formula}s دے کثیر سیٹ نیں جیہڑی اسیں ورت رہے آں۔ رواج مطابق ایہ وی اے کہ $\tuple{\Gamma, \Delta}$ لکھن دی تھاں ثبوت دے نظریے دے ماہر دوہاں جزواں نوں سیکوئنٹ دی نشانی نال وکھ کردے نیں۔ اسیں~$\Sequent$ ورتدے آں۔

\begin{defn}[سیکوئنٹ]
اک \emph{سیکوئنٹ} ایس شکل دی عبارت اے:
\[
\Gamma \Sequent \Delta
\]
جتھے $\Gamma$ تے $\Delta$، !!{formula}s دے محدود (شاید خالی) کثیر سیٹ نیں۔ $\Gamma$ نوں \emph{کھبّا پاسا} تے $\Delta$ نوں \emph{سجّا پاسا} آکھدے آں۔
\end{defn}

منّ لو $!G$، اوہناں !!{formula}s دا عطف اے جیہڑے~$\Gamma$ وچ نیں (خالی عطف نوں~$\ltrue$ منّ کے)، تے $!D$، اوہناں !!{formula}s دی فصل اے جیہڑے~$\Delta$ وچ نیں (خالی فصل نوں~$\lfalse$ منّ کے)۔ تاں سیکوئنٹ $\Gamma \Sequent \Delta$ کسے ساخت~$\Struct{M}$ وچ اوہدوں تے صرف اوہدوں سچا اے جدوں $!G \lif !D$ سچا اے۔ کلاسیکی منطق وچ ایہ اوہو گل اے کہ یا اک !!{formula}s وچوں جیہڑا~$\Gamma$ وچ اے جھوٹھا اے، یا اک !!{formula}s وچوں جیہڑا~$\Delta$ وچ اے سچا اے۔

جے $\Gamma$، !!{formula}s دا کثیر سیٹ اے، تاں اسیں $\Gamma, !A$ (یا $!A, \Gamma$) لکھدے آں، $!A$ نوں~$\Gamma$ وچ شامل کرن دے حاصل لئی۔ جے $ \Delta$ وی !!{formula}s دا کثیر سیٹ اے، تاں $\Gamma, \Delta$ دوہاں لڑیاں دا کثیر سیٹ اتحاد اے۔ جے $\Gamma$ وچ $!A$ دی \emph{کثرت}~n اے (یعنی اوہدے وچ $n$ نقلاں~$!A$ دیاں نیں) تے $ \Delta$ وچ کثرت~$m$ اے، تاں $\Gamma, \Delta$ وچ $!A$ دی کثرت~$n+m$ اے۔ جے سیکوئنٹ دے اک پاسے دا کثیر سیٹ خالی اے، تاں اسیں عام طور اُتے اوس پاسے خالی تھاں چھڈّ دیندے آں۔ مثال لئی $\Sequent !A$ اوہ سیکوئنٹ اے جیہدا کھبّا پاسا خالی اے تے سجّے پاسے $!A$ دی کثرت~$1$ اے۔

سیکوئنٹ حساب وچ !!^a{proof} سیکوئنٹاں دا درخت اے۔ اوہدے سبھ توں اُتلے (یعنی شروع دے) سیکوئنٹ متعلقہ نظام دے مسلّمے نیں، تے جے کوئی سیکوئنٹ اک یا دو ہور سیکوئنٹاں دے تھلّے اے، تاں اوہ نظام دے کسے استنباطی قاعدے نال ٹھیک طرح نکلنا چاہیدا اے۔ اسیں پہلاں کلاسیکی منطق لئی دو وکھرے سیکوئنٹ نظام، $\Log{G1c}$ تے~$\Log{G3c}$، ویکھاں گے۔ اوہناں دے مسلّمے تے قاعدے \cref{pt:seq:int:tab:G1c,pt:seq:int:tab:G3c} وچ دِتّے نیں۔ \oliflabeldef{fol:seq::chap}{ (گینٹسن دا اصل نظام \Log{LK}، \olref[fol][seq][]{chap} وچ بیان کیتا گیا اے۔)}{}

\subfile{rules-G1c}
\subfile{rules-G3c}

نظاماں وچ باریک فرق نیں تے اوہناں کر کے ثبوت دے نظریے نال متعلق اوہناں دیاں خاصیتاں وکھریاں نیں۔ \Log{G1c} دے مسلّمے $!A
\Sequent !A$ نیں، جیہناں وچ ہور !!{formula}s دی کھُل نہیں۔ \Log{G3c} دے مسلّمے $!C, \Gamma \Sequent \Delta, !C$ دی شکل دے نیں، جیہناں وچ «پاسے دے» کوئی وی !!{formula}s، $\Gamma$ تے~$\Delta$ وچ، ہو سکدے نیں، پر !!{formula}~$!C$ جیہڑا دوہاں پاسیاں اُتے آؤندا اے \emph{جوہری} ہونا چاہیدا اے۔ \Log{G1c} وچ $\LeftR{\land}$ تے $\RightR{\lor}$ قاعدیاں دیاں دو دو شکلاں نیں، جدکہ \Log{G3c} وچ اک اک اے۔ \Log{G1c} وچ وادھو «ساختی قاعدے» وی نیں، جنہاں نوں «سکڑاؤ» (\LeftR{\Contraction}، \RightR{\Contraction}) تے «کمزور کرنا» (\LeftR{\Weakening}، \RightR{\Weakening}) آکھدے آں۔

\Log{G1c} تے \Log{G3c} نظام کلاسیکی منطق لئی صحت مند تے کامل نیں۔ دوجے لفظاں وچ، ہر سیکوئنٹ جیہڑا ایہناں نظاماں وچ !!{provable} اے، ہر !!{structure} وچ سچا اے، تے ہر سیکوئنٹ جیہڑا ہر !!{structure} وچ سچا اے، اوہدا !!a{proof} اے۔ سو ایہ سوال کہ \emph{کی} کسے سیکوئنٹ دا !!a{proof} اے، منطق دے ہور طریقیاں نال وی حل ہو سکدا اے (جیویں نمونیاں دے نظریے دے طریقیاں نال)۔ تے کیونکہ دونے نظام عین اوہ سیکوئنٹ !!{prove} کردے نیں جیہڑے ہر !!{structure} وچ سچے نیں، ایس لئی خاص کر اوہ اوہی سیکوئنٹ ثابت کردے نیں۔

پر ثبوت دے نظریے دا واسطہ صرف ایس سوال نال نہیں کہ \emph{کی} کوئی گل ثابت ہو سکدی اے، بلکہ ایہ وی کہ \emph{کِنج}۔ ثبوت دے نظریے دے ماہر ایہو جیہے سوال ویکھدے نیں، «کی کسے خاص قاعدے نوں ہمیشہ ٹالیا جا سکدا اے؟»، «کی ایہ قاعدہ نظام وچ شامل کیتا جا سکدا اے بغیر ایہدے کہ نویں سیکوئنٹ ثابت ہون جوگ ہو جان؟»، یا «کی اک نظام دے !!{proof}s دوجے دے !!{proof}s وچ بدلے جا سکدے نیں؟» جے ایہو جیہے سوال دا جواب «ہاں» اے، تاں ایس گل دا ثبوت اکثر !!{proof}s دی گنجل اُتے استقرا نال، یا اوہو گل دوجے لفظاں وچ، !!{proof}s دی بناوٹ اُتے استقرا نال ملدا اے۔ ایس نال گل دا ثبوت ملدا اے (مثلاً جے \Log{G1c} اک سیکوئنٹ ثابت کردا اے تاں \Log{G3c} وی اوہی سیکوئنٹ ثابت کردا اے)، تے اک طریقہ وی (مثلاً !!{proof}s،~\Log{G1c} دے، نوں !!{proof}s،~\Log{G3c} دے، وچ بدلّن دا)۔

ثبوت دے نظریے دا اک مرکزی نتیجہ \emph{قطع ہٹاؤن دا قضیہ} اے۔ ایہ وکھاؤندا اے کہ اسیں قطع دا قاعدہ
\[
\Axiom$\Gamma \fCenter \Delta, !A$
\Axiom$!A, \Pi \fCenter \Lambda$
\RightLabel{\Cut}
\BinaryInf$\Gamma, \Pi \fCenter \Delta, \Lambda$
\DisplayProof
\]
سیکوئنٹ نظاماں وچ شامل کر سکدے آں، بغیر ایہدے کہ !!{provable} سیکوئنٹ وکھرے ہو جان۔ نالے ثبوت اک طریقہ دیندا اے جیہڑا کوئی وی !!{proof}، جیہڑا \Log{G1c} (یا \Log{G3c}) دے قاعدے تے نال \Cut{} ورتدا اے، ایہو جیہے ثبوت وچ بدل دیندا اے جیہڑا صرف \Log{G1c} (یا \Log{G3c}) دے قاعدے ورتدا اے۔ دوجے لفظاں وچ، ایہ طریقہ \Cut{} قاعدے نوں اوہناں !!{proof}s وچوں «ہٹا دیندا» اے جیہڑے اوہنوں ورتدے نیں۔

\end{document}
